\documentclass[landscape,a4paper]{article}

\usepackage[landscape, top=0.3cm, bottom=0.3cm, left=0.3cm, right=0.3cm]{geometry}
\usepackage{multicol}

% 2. 字体编码和基础数学包
\usepackage[T1]{fontenc}
\usepackage{amsmath,amsthm,amsfonts,amssymb}
\usepackage{mathtools}
\usepackage{derivative}
\usepackage{upgreek}
\usepackage[condensed,math]{iwona}
\def\Ftrans{\raisebox{4pt}{$\xleftrightarrow{\mathcal{F}}$}}
\def\FTrans{\xleftrightarrow{\mathcal{F}}}
\def\Ltrans{\raisebox{4pt}{$\xleftrightarrow{\mathcal{L}}$}}
\usepackage{enumitem}
\setlist[itemize,2]{label=\fontsize{3pt}{3.5pt}\selectfont$\square$}
\usepackage{tikz}
\usetikzlibrary{positioning, calc}
\setlength{\parindent}{0pt}
\setlength{\parskip}{0pt}
\setlength{\columnsep}{3pt}
\usepackage{xcolor}
\usepackage{soul}

% Set the highlight color to a bright yellow
\sethlcolor{yellow} 

% Save the original \emph command just in case
\let\oldemph\emph 

% Redefine \emph to use the soul highlight command
% (Optional: Add \textit{#1} inside \hl if you want to keep the italics as well)
\renewcommand{\emph}[1]{\hl{#1}}
\usepackage{enumitem}
\setlist{nosep, leftmargin=0.4em, labelsep=0.2em} 

\usepackage[most]{tcolorbox}

\newtcolorbox{card}[1]{
	enhanced,
	sharp corners,
	breakable,
	colback=white,
	colbacktitle=gray!10,
	coltitle=black,
	fonttitle=\bfseries\fontseries{bc}\selectfont\scriptsize, 
	title={#1},
	frame style={
		draw=black!70,
		line width=0.5pt,
		dashed,
		dash pattern=on 2pt off 1pt
	},
	left=1pt, right=1pt, top=1pt, bottom=1pt,
	toptitle=1pt, bottomtitle=1pt,
	boxrule=0pt,
	before skip=0pt,
	after skip=0pt,
	parbox=false,
}

\usepackage{listings}
\lstset{
	basicstyle=\ttfamily\fontsize{6pt}{6pt}\selectfont,
	breaklines=true,
	frame=none,
	aboveskip=0pt, belowskip=0pt
}

% -----------------------------------------------------------------------
% 正文
% -----------------------------------------------------------------------
\newcommand{\ii}{\mathrm{i}}
\newcommand{\e}{\mathrm{e}}
\newcommand*{\oo}[1]{\left(#1\right)}
\newcommand*{\oc}[1]{\left(#1\right]}
\newcommand*{\cc}[1]{\left[#1\right]}
\newcommand*{\co}[1]{\left[#1\right)}
\newcommand{\C}[2]{\mathrm{C}_{#1}^{#2}}
\newcommand{\PP}{\mathbf{P}}
\newcommand{\EE}{\mathbf{E}}
\newcommand{\var}{\operatorname{var}}
\begin{document}
	
	\fontsize{7pt}{8pt}\selectfont 

	
	\begin{multicols*}{4}
		\begin{card}{Combination Numbers}
			\begin{itemize}
				\item Pascal's Rule: $\C{n}{k}=\C{n-1}{k}+\C{n-1}{k-1}$.
				\item \emph{Absorption Identity: $k\C{n}{k}=n\C{n-1}{k-1}$.}
				\item Trinomial Revision: $\C{n}{m}\C{m}{k}=\C{n}{k}\C{n-k}{m-k}$.
				\item Row Sum: $\sum_{k=0}{n}\C{n}{k}=2^n$.
				\item Even-Odd Sums: $\sum_{k\ge 0} \C{n}{2k}=\sum_{k\ge 0}\C{n}{2k+1}=2^{n-1}$.
				\item Weighted Sums:
				\begin{itemize} 
					\item $\sum_{k=0}^{n}k\C{n}{k}=n2^{n-1}$.
				\item $\sum_{k=0}^{n}k^2\C{n}{k}=n(n+1)2^{n-2}$.
				\item  $\sum_{k=0}^{n}k^3\C{n}{k}=n^2(n+3)2^{n-3}$.
				\item  $\sum_{k=0}^m(-1)^k\C{n}{k}=(-1)^m\C{n-1}{m}$.
				\end{itemize}
				\item Hockey-Stick Identity: $\sum_{i=k}^n\C{i}{k}=\C{n+1}{k+1}$.
				\item Diagonal Sum: $\sum_{k=0}^m \C{n+k}{k}=\C{n+m+1}{m}$.
				\item Vandermonde's Identity: $\sum_{k=0}^r \C{n}{k}\C{m}{r-k}=\C{n+m}{r}$.
			\end{itemize}
		\end{card}
		\begin{card}{Probability Space}
			\begin{itemize}
				\item \textbf{Sample Space $\Omega$}: the set of all possible outcomes.
				\item \textbf{$\upsigma$-algebra}: the collection of events.
				\begin{itemize}
					\item $\emptyset\in\mathcal F$,
					\item If $A\in\mathcal F$, then $A^\complement\in \mathcal F$,
					\item If $A_i\in\mathcal F$, then $\bigcup_i A_i\in\mathcal F$.
				\end{itemize}
				\item \textbf{Borel field}: all intervals and countable union of them.
				\item \textbf{Probability Measure $\PP$}: a function $\PP \colon \mathcal F\to [0,1]$.
				\begin{itemize}
					\item Nonnegativity: $\PP(A)\ge 0,\forall A\in\mathcal F$,
					\item Normalization: $\PP(\emptyset)=0,\PP(\Omega)=1$,
					\item Countable additivity: If $A_i\in\mathcal F$ are \textit{disjoint}, $\PP\left(\bigcup_i A_i\right)=\sum_i\PP(A_i)$.
				\end{itemize}
			\end{itemize}
		\end{card}
		\begin{card}{Conditioning and Independence}
			\makebox[0.9\textwidth][l]{\textbf{\textit{Conditioning}}\hrulefill}
			\begin{itemize}
				\item \textbf{Definition}: For $\PP(B)\neq 0$, $\PP(A|B)=\PP(A\cap B)/\PP(B)$.
				\begin{itemize}
					\item Nonnegativity: $\PP(A|B)\ge 0,\forall A\in\mathcal F$,
					\item Normalization: $\PP(\Omega|B)=1$,
					\item Countable additivity: If $A_i\in\mathcal F$ are \textit{disjoint}, $\PP\left(\bigcup_i A_i|B\right)=\sum_i\PP(A_i|B)$.
				\end{itemize}
				\item \textbf{Multiplication Rule}: $\PP(A\cap B\cap C)=\PP(A)\cdot\PP(B|A)\cdot\PP(C|A\cap B)$.
				\item \textbf{Total Probability Theorem}: For a partition of sample space $A_1, A_2,\dots, A_n$, assuming $\PP(A_i)>0$ for all $i$,
				\[
				\PP(B)=\sum_{i=1}^n \PP(A_i)\PP(B|A_i).
				\]
				\item \textbf{Bayes's Rule}: $\PP(A_i|B)=\dfrac{\PP(A_i)\PP(B|A_i)}{\sum_{j=1}^n\PP(A_j)\PP(B|A_j)}$.
			\end{itemize}
			\makebox[0.9\textwidth][l]{\textbf{\textit{Independence}}\hrulefill}
			\begin{itemize}
			\item \textbf{Definition}: Event $A$ and $B$ are called \textbf{independent} if $\PP(A\cap B)=\PP(A)\PP(B)$.
			\item For $\PP(B)>0$, $\PP(A|B)=\PP(A)$.
			\item \textbf{Conditional Independence}: $\PP(A\cap B|C)=\PP(A|C)\PP(B|C)$.
			\item \textbf{Collective Independence}: Events $A_1,A_2,\dots, A_n$ are called independent if
			\[
			\PP(A_i\cap A_j\cap \dots\cap A_q)=\PP(A_i)\PP(A_j)\dots\PP(A_q),
			\]
			for any distince indices $i,j,\dots,q$ chosen from $\left\{i,j,\dots,q\right\}$.
			\item \emph{\textbf{Inclusion-Exclusion Principle}:}
			$\PP\left(\bigcup_{i=1}^n A_i\right) = \sum \PP(A_i) - \sum_{i<j} \PP(A_i \cap A_j) + \dots + (-1)^{n-1} \PP(A_1 \cap \dots \cap A_n)$.
			\begin{itemize}
				\item \textbf{Hat Problem (Derangement)}: 
				\begin{itemize}
					\item Derangement Number:\\ $D_n=n!\sum_{i=0}^n(-1)^i/i!$.\\
					$D_n=nD_{n-1}+(-1)^n$.\\
					$\lim_{n\to\infty} D_n/n!=1/\e$.
					\item Prob. of at least one match is $\approx 1 - 1/\e$ as $n\to\infty$.
				\end{itemize}
				\item \textbf{Symmetric Difference}: $\PP(A \triangle B) = \PP(A) + \PP(B) - 2\PP(A \cap B)$.
			\end{itemize}
			\end{itemize}
		\end{card}
		\begin{card}{Discrete Random Variable}
			\begin{itemize}
				\item \textbf{Random Variable}: Given a probability space $(\Omega,\mathcal F,\PP)$, a \textbf{random variable} is a \textit{funtion} $X\colon \Omega\to R$ such that $X^{-1}(A)\in\mathcal F$ for each $A\in\mathcal B(\mathbb R)$.
				\item \textbf{PMF}:\\ $p_X(x)=\PP(X=x)$.
				\item \textbf{Expectation}:\\ $\EE[X]=\sum_ix_ip_X(x_i)$.
				\item \textbf{Variance}:\\ $\var[X]=\EE\left[\left(X-\EE(X)\right)^2\right]=\EE[X^2]-\left(\EE[X]\right)^2$.
				\item \emph{$\EE[\alpha X+ \beta]=\alpha\EE[X]+\beta,\EE[X+Y]=\EE[X]+\EE[Y].$}
				\item $\var[\alpha X+\beta]=\alpha^2\var[X]$.
				\item \begin{itemize}
					\item \textbf{Geometric RV}:
					\begin{itemize}
						\item $p_X(k)=(1-p)^{k-1}p$,
						\item $\EE[X]=1/p$,
						\item $\EE[X^2]=1/p$,
						\item $\var[X]=(1-p)/p^2$.
						\item \emph{\textbf{Memorylessness}:}\\ $p_X(k|X>n)=(1-p)^{k-n-1}p;\EE[X|X > n]=1/p + n$.
					\end{itemize}
					\item \textbf{Bernoilli RV}:
					\begin{itemize}
						\item $p_X(k)=\begin{cases}p,&k=1,\\1-p,&k=0,\end{cases}$,
						\item $\EE[X]=p$,
						\item $\EE[X^2]=p$,
						\item $\var[X]=p(1-p)$.
					\end{itemize}
					\item \textbf{Discrete Uniform RV}:
					\begin{itemize}
						\item $p_X(k)=1/(b-a+1)$, if $k=a,a+1,\dots,b$,
						\item $\EE[X]=(a+b)/2$,
						\item $\var[X]=(b-a)(b-a+2)/12$.
					\end{itemize}
					\item \textbf{Binomial RV}:
					\begin{itemize}
						\item $p_X(k)=\C{n}{k}p^k(1-p)^{n-k}$,
						\item $\EE[X]=np$,
						\item $\EE[X^2]=np(1-p+np)$,
						\item $\var[X]=np(1-p)$.
					\end{itemize}
					\item \textbf{Poisson RV}: \emph{$\lambda=np$ as $n\to\infty$ in binomial RV.}
					\begin{itemize}
						\item $p_X(k)=\e^{-\lambda}\lambda^k/k!,k\in\mathbb N$,
						\item $\EE[X]=\lambda$,
						\item $\EE[X^2]=\lambda+\lambda^2$,
						\item $\var[X]=\lambda$,
						\item \emph{$X_1+X_2\sim\text{Poisson}(\lambda_1+\lambda_2)$.}
					\end{itemize}
				\end{itemize}
				\end{itemize}
			\end{card}
			\begin{card}{Multiple Discrete RV}
				\begin{itemize}
					\item \textbf{Conditioning}:
					\begin{itemize}
						\item \textbf{Conditional PMF}:\\ $p_{X|A}(x)=\PP(X=x\cup A)/\PP(A),p_{X|Y}(x|y)=p_{X,Y}(x,y)/p_Y(y)$.
						\item \textbf{Conditional Expectation}:\\ $\EE[X|A]=\sum_x xp_{X|A}(x)$.
						\item \emph{\textbf{Total Expectation Theorem}:}\\ $\EE[X]=\PP(A_1)\EE[X|A_1]+\dots +\PP(A_n)\EE[X|A_n]$.
						\item \textbf{Marginal Distribution Function}:\\ $p_{X,Y}(x,y)=\sum_z p_{X,Y,Z}(x,y,z);p_X(x)=\sum_y p_{X,Y}(x,y)$.
					\end{itemize}
					\item \textbf{Independence}:\\ $\forall x,y\colon p_{X,Y}(x,y)=p_X(x)p_Y(y)$.
					\item For Independent RVs:
					\begin{itemize}
						\item $f(X)$ and $g(Y)$ are also independent.
						\item \emph{$\EE[XY]=\EE[X]\EE[Y]$.}
						\item \emph{$\var(X+Y)=\var(X)+\var(Y), \var(X-Y)=\var(X)+\var(Y)$.}
					\end{itemize}
					\item \textbf{Conditional Independence}:\\ $\forall x,y\colon p_{X,Y}(x,y|A)=p_X(x|A)p_Y(y|a)$.
				\end{itemize}
		\end{card}
		\begin{card}{Continuous RV}
			\begin{itemize}
				\item \textbf{PDF}:\\ $\PP(a\le X\le b)=\int_a^bf_X(x)\,\mathrm dx$;\\
				$\PP(x\le X\le x+\delta)\approx f_X(x)\delta$.
				\item \textbf{Expectation}:\\ $\EE[X]=\int_{-\infty}^{+\infty} xf_X(x)\,\mathrm dx$.
				\item \textbf{Variance}:\\ $\var(X)=\int_{-\infty}^{+\infty}\left(x-\EE[X]\right)^2f_X(x)\,\mathrm dx=\EE[X^2]-\left(\EE[X]\right)^2$.
				\item \begin{itemize}
					\item \textbf{Uniform RV}:
					\begin{itemize}
						\item $f_X(x)=1/(b-a),a\le x\le b$,
						\item $\EE[X]=(a+b)/2$,
						\item $\EE[X^2]=(a^2+ab+b^2)/3$,
						\item $\var[X]=(b-a)^2/12$.
					\end{itemize}
					\item \textbf{Exponential RV}:
					\begin{itemize}
						\item $f_X(x)=\lambda\e^{-\lambda x},x\ge 0$,
						\item $F_X(x)=\e^{-\lambda x}$.
						\item $\EE[X]=1/\lambda$,
						\item $\EE[X^2]=2/\lambda^2$,
						\item $\var[X]=1/\lambda^2$.
						\item \emph{$\PP(X_1<X_2)=\lambda_1/(\lambda_1+\lambda_2)$.}
						\item Memorylessness:\\
						$f_X(x|X>a)=\lambda \e^{-\lambda(x-a)}, \PP(X\ge c+x|X\ge c)=\PP(X\ge x)=\e^{-\lambda x};$\\$\EE[X|X>c]=1/\lambda+c$.
					\end{itemize}
					\item \textbf{Gaussian RV} ($X\sim N(\mu,\sigma^2)$):
					\begin{itemize}
						\item \[f_X(x)=\frac{1}{\sigma\sqrt{2\uppi}}\e^{-(x-\mu)^2/2\sigma^2},\]
						\item $F_\Lambda(\lambda)=\varPhi(\lambda),\Lambda\sim N(0,1)$.
						\item $\EE[X]=\mu$,
						\item $\EE[X^2]=\mu^2+\sigma^2$,
						\item $\EE[X^4]=3\sigma^4+6\sigma^2\mu^2+\mu^4$,
						\item $\var[X]=\sigma^2$.
						\item $Y=\alpha X+\beta\sim N(\alpha\mu+\beta,\alpha^2\sigma^2)$.
					\end{itemize}
				\end{itemize}
				\item \textbf{CDF}:\begin{itemize}
					\item $F_X(x)=\int_{-\infty}^{x_+} f_X(t)\,\mathrm dt$.
					\item $F_X(x)$ uncontinuous $\Rightarrow$ $X$ has discrete part.
					\item $\PP(a<X\le b)=F_X(b)-F_X(a)$.
					\item $\EE[X]=\int_0^{+\infty}(1-F_X(x))\,\mathrm dx-\int_{-\infty}^0F_X(x)\,\mathrm dx$.
				\end{itemize}
				
				\item \emph{\textbf{Min and Max of i.i.d. RVs} ($X_1, \dots, X_n$):}
				\begin{itemize}
					\item \textbf{Max}: $Z = \max(X_i) \implies F_Z(z) = [F_X(z)]^n \implies f_Z(z) = n[F_X(z)]^{n-1} f_X(z)$.
					\item \textbf{Min}: $W = \min(X_i) \implies F_W(w) = 1 - [1 - F_X(w)]^n$.
					\item \textit{Shortcut}: If $X_i \sim \operatorname{Exp}(\lambda_i)$, then $\min(X_1, \dots, X_n) \sim \operatorname{Exp}(\sum \lambda_i)$.
				\end{itemize}
				
				\item \textbf{2D Uniform \& Geometric Probability}:
				\begin{itemize}
					\item $f_{X,Y}(x,y) = 1/\text{Area}(S)$ for $(x,y) \in S$.
					\item \textbf{Broken Stick / Triangle Inequality}: Break a line into 3 pieces. They form a triangle iff. all pieces are shorter than $1/2$. Prob. is $1/4$.
				\end{itemize}
			\end{itemize}
		\end{card}
		\begin{card}{Multiple Continuous RV}
			\textbf{Contidioning}:
			\begin{itemize}
				\item \textbf{Conditional PDF}: $f_{X|Y}(x|y)=f_{X,Y}(x,y)/f_{Y}(y)$.
				\item \textbf{Conditional Expectation}: $\EE[X|Y=y]=\int_{-\infty}^{+\infty}xf_{X|Y}(x|y)\,\mathrm dx$.
				\item \emph{\textbf{Total Expectation Theorem}: $\EE[X]=\int_{-\infty}^{+\infty}\EE[X|Y=y]f_Y(y)\,\mathrm dy$.}
				\item \textbf{Marginal Distribution Function}: $f_X(x)=\int_{-\infty}^{+\infty} f_{X,Y}(x,y)f_Y(y)\,\mathrm dy$.
			\end{itemize}
			\textbf{Independence}:
			\begin{itemize}
				\item $\forall x,y\colon f_{X,Y}(x,y)=f_X(x)f_Y(y)$.
				\item For independent RVs:
				\begin{itemize}
					\item $f(X)$ and $g(Y)$ are also independent,
					\item $\EE[XY]=\EE[X]\EE[Y]$,
					\item $\var(X+Y)=\var(X)+\var(Y)$.
					\item \textbf{Convolution}: If $X$ and $Y$ are independent and $Z=X+Y$, $f_Z(z)=f_X(x)\ast f_Y(y)=\int_{-\infty}^{+\infty}f_X(x)f_Y(z-x)\,\mathrm dx$.
				\end{itemize}
			\end{itemize}
			\textbf{Joint / Multiple CDF}: $F_{X,Y}(x,y) = \PP(X \le x, Y \le y)$.
			\begin{itemize}
				\item \textbf{Marginals}: $F_X(x) = F_{X,Y}(x, \infty)$, $F_Y(y) = F_{X,Y}(\infty, y)$.
				\item \textbf{From/To PDF}:\\$F_{X,Y}(x,y) = \int_{-\infty}^{x} \int_{-\infty}^{y} f_{X,Y}(u,v)\,\mathrm dv\,\mathrm du $\\$\implies f_{X,Y}(x,y) = \frac{\partial^2}{\partial x \partial y} F_{X,Y}(x,y)$.
				\item \textbf{Independence}: $X \perp Y \iff F_{X,Y}(x,y) = F_X(x)F_Y(y)$.
			\end{itemize}
			
			\textbf{Conditional CDF}: 
			\begin{itemize}
				\item \textbf{Given Event $A$}: $F_{X|A}(x) = \PP(X \le x \mid A) = \frac{\PP(\{X \le x\} \cap A)}{\PP(A)}$.
				\item \textbf{Given $Y=y$}: $F_{X|Y}(x|y) = \int_{-\infty}^{x} f_{X|Y}(t|y) \,\mathrm dt$.
				\item \textbf{Derivative}: $f_{X|Y}(x|y) = \frac{\partial}{\partial x} F_{X|Y}(x|y)$.
				\item \textbf{Total Prob}: $F_X(x) = \sum \PP(A_i) F_{X|A_i}(x)$.
			\end{itemize}
			\textbf{Continuous Bayes' Rule}:
			\[
			f_{X | Y}(x | y)=\frac{f_{Y | X}(y | x) f_X(x)}{f_Y(y)}=\frac{f_{Y | X}(y | x) f_X(x)}{\int_{-\infty}^{\infty} f_{Y | X}(y | t) f_X(t) d t},
			\]
			\[\PP(A|Y=y) = \frac{\PP(A) f_{Y|A}(y)}{f_Y(y)}. (\text{Mixed Bayes})\]
		\end{card}
		\begin{card}{Intergration Formula}
			\begin{itemize}
				\item $\int x\e^{ax}\,\mathrm dx=\frac{1}{a^2}(ax-1)\e^{ax},\int x^2\e^{ax}\,\mathrm dx=\frac{a^2x^2-2ax+2}{a^3}\e^{ax},\int x^3\e^{ax}\,\mathrm dx=\frac{a^3x^3-3a^2x^2+6ax-6}{a^3}\e^{ax},\int x^n\e^{ax}\,\mathrm dx=\frac{1}{a}x^n\e^{ax}-\frac{n}{a}\int x^{n-1}\e^{ax}\,\mathrm dx.$
				\item $\int_0^{\infty}x^n\e^{-ax}\,\mathrm dx=n!/a^{n+1}.$
				\item $\int\sqrt{a^2-x^2}\,\mathrm dx=\frac{x}{2}\sqrt{a^2-x^2}+\frac{a^2}{2}\arcsin\frac{x}{a}.$
				\item $\int \ln x\,\mathrm dx=x\ln x-x,\int x^n\ln x\,\mathrm dx=\frac{1}{n+1}x^{n+1}\left(\ln x-\frac{1}{n+1}\right).$
				\item $\int_{-\infty}^{+\infty} \exp\left(-(x-\lambda)^2/\mu^2\right)\,\mathrm dx=\mu\sqrt{\uppi}$.
				\item $\int_{-\infty}^{+\infty}x^2\exp\left(-x^2/\mu^2\right)\,\mathrm dx=\frac{\mu}{2}\sqrt{\uppi}.$
			\end{itemize}
		\end{card}
	\end{multicols*}
\end{document}